% docs/style/uni-packages.tex -- local copy of tex-skripten/formatierungen/uni-packages.tex
% Copied 2026-10-12 so that the repository compiles stand-alone (GitHub).
% Local version: 1.0.0 -- language switched to English (project language), see markers "interferosim"

%====================================================
%  UNI-Dokumente – Zentrale Präambel
%  (Editorbreite 59 Zeichen, PDF 50 %)
%====================================================

%----------------------------------------------------
% 1. Dokumentbasis & Sprache
%----------------------------------------------------
%\usepackage[naustrian]{babel}    % deutsche Verzeichnisse
\usepackage[english]{babel}      % interferosim: project language is English
\usepackage{textcomp}            % Zusatzsymbole (z. B. €)
\usepackage{xspace}

%----------------------------------------------------
% 2. Layout, Kopf- und Fußzeilen
%----------------------------------------------------
\usepackage[
headsepline=on,
footsepline=on,
automark
]{scrlayer-scrpage}

\usepackage{standalone}
\usepackage{lscape}
\usepackage{rotating}
\newcount\Sourcenummer\Sourcenummer=0
\usepackage{afterpage}
\usepackage{flafter}
\usepackage[auto]{chappg}
\usepackage{changepage}          % adjustwidth

%----------------------------------------------------
% 3. Schriften
%----------------------------------------------------
\usepackage{helvet}
\usepackage{courier}

%----------------------------------------------------
% 4. Farben & Boxen
%----------------------------------------------------
\usepackage[table,x11names,dvipsnames]{xcolor}
\usepackage{tcolorbox}
\usepackage{framed}

%----------------------------------------------------
% 5. Listen & Aufzählungen
%----------------------------------------------------
\usepackage{paralist}
\usepackage{enumerate}
\usepackage{enumitem}
\usepackage{mdwlist}
\usepackage{multicol}
\usepackage{parallel}

%----------------------------------------------------
% 6. Code & Verbatim
%----------------------------------------------------
\usepackage{listings}[showspaces=false]
\usepackage{verbatim}
\usepackage{spverbatim}

%----------------------------------------------------
% 7. Grafiken & Gleitobjekte
%----------------------------------------------------
\usepackage{graphicx}
\usepackage{caption}
\usepackage{pdflscape}
\usepackage{pdfpages}
\usepackage{thumbpdf}

%----------------------------------------------------
% 8. Mathematik & Physik
%----------------------------------------------------
\usepackage{mathtools}
\usepackage{amsmath}
\usepackage{amssymb}
\usepackage{mathrsfs}
\usepackage{trfsigns}
\usepackage{color,cancel}
\usepackage{xparse}
\usepackage{physics}
\usepackage[version=4]{mhchem}
\usepackage[autostyle]{csquotes} % Behebt die Babel-Warnung

\usepackage{siunitx}
%\sisetup{locale = DE}
\sisetup{locale = UK}             % interferosim: decimal point

%----------------------------------------------------
% 9. Tabellen
%----------------------------------------------------
\usepackage{array}       % für >{\raggedright\arraybackslash}-Spaltentypen (z.B. Formelverzeichnis)
\usepackage{booktabs}
\usepackage{multirow}
\usepackage{diagbox}
% WICHTIG: longtable muss VOR arydshln geladen werden, damit dessen
% interner longtable-Kompatibilitätspatch (\LT@array) aktiv wird –
% sonst Fehler "Undefined control sequence \adl@..." im
% Formelzeichenverzeichnis (das intern longtable verwendet).
\usepackage{longtable}
\usepackage{arydshln}
\usepackage{csvsimple}
\usepackage{tabularray}   % Ersatz/Erweiterung für longtable

%----------------------------------------------------
% 10. PDF, Hyperlinks & Navigation
%----------------------------------------------------
\usepackage[
colorlinks=true,
plainpages=false, 
pdfpagelabels,
hypertexnames=false % Wichtig: Behebt die "destination with same identifier" Warnung bei Gleichungen
]{hyperref}
\usepackage{bookmark}

%----------------------------------------------------
% 11. Indizes, Glossare & Bibliografie
%----------------------------------------------------
\usepackage{imakeidx}

\usepackage[
automake,
nonumberlist,
acronym,
nopostdot,
toc=false,
nomain
]{glossaries-extra}

\usepackage[
backend=biber,
style=numeric,
sorting=ynt
]{biblatex}

%----------------------------------------------------
% 12. Diagramme & Zeichnungen
%----------------------------------------------------
\usepackage{tikz}
% WICHTIG: arrows.meta muss geladen sein für 'Stealth'/Latex-Pfeilspitzen;
% alle vom Skriptum benötigten Libraries in einer Zeile vereinheitlicht:
\usetikzlibrary{
	positioning,
	calc,
	arrows.meta,
	angles,
	quotes,
	shapes.geometric,
	patterns,
	backgrounds
}
\usepackage{pgfplots}
\pgfplotsset{compat=1.18}

%----------------------------------------------------
% 13. Spezial-/Legacy-Pakete (PSTricks)
%----------------------------------------------------
\usepackage{pst-pdf}
% \usepackage{pst-circ} % optional
